\documentclass[10pt,a4paper]{amsart}
\usepackage[margin=19mm]{geometry}
\usepackage{amsmath,amssymb,mathtools}
\usepackage[scr=rsfs,cal=euler]{mathalfa}
\usepackage{xfrac}
\usepackage[unicode,hidelinks,bookmarks=false]{hyperref}
\newcommand{\ie}{i.e.\ }
\newcommand{\cf}{cf.\ }
\newcommand{\resp}{respectively\ }
\newcommand{\Subsection}{\subsection}
\providecommand{\nobreakdash}{\nobreak\hbox{-}}
\setlength{\emergencystretch}{2em}
\multlinegap0pt
\usepackage{amssymb}
\newtheorem{theorem}{Theorem}
\newtheorem{corollary}{Corollary}
\usepackage{cleveref}
\crefname{theorem}{Theorem}{Theorems}
\crefname{corollary}{Corollary}{Corollarys}
\title{General formulas for a class of Euler sums}
\author{David H Bailey, Ross McPhedran, Bruno Salvy}
\date{Source: arXiv:2604.02384v1 (CC BY 4.0)}
\begin{document}
\maketitle

\noindent\emph{Source and licence.} General formulas for a class of Euler sums. Version arXiv:2604.02384v1 (CC BY 4.0), distributed by arXiv under the Creative Commons Attribution 4.0 International licence, http://creativecommons.org/licenses/by/4.0/, as recorded in the arXiv licence metadata for this exact version and shown on the abstract page https://arxiv.org/abs/2604.02384v1. Full licence notice: \texttt{licenses/arxiv-2604.02384-CC-BY-4.0.txt}.\par\smallskip

\noindent\textit{Editorial scope.} Counted unit: the article's corollary cor:cor2 (source0.tex lines 420--426). The article's complete original proof of it is delivered with it (source0.tex lines 427--436, one \texttt{proof} environment of 10 lines). No mathematics is written, repaired or re-derived: the statement and the proof are the article's own lines, in the article's own notation, and no retained line is substituted or re-flowed.  Retained uncounted as the source-local context the proof needs: lines 345--351; lines 403--415.  Nothing was omitted from any retained span, so the package carries no omission record.  No retained line contains a citation, so the wrapper carries no bibliography and no bibliography entry.  No retained line contains a figure environment, an image-inclusion command or drawing code, so no image is delivered and the package declares no asset dependency.  Editorial formatting only, in addition: an AMS \texttt{amsart} preamble that restates the article's own declarations and macros used by the retained lines, namely the environments \texttt{theorem}, \texttt{corollary} on separate counters, together with the standard packages those lines need.  The retained environments print in this document as Theorem 1, Corollary 1 in order of appearance, whereas the article prints the same items with the numbering of its own full document.  The complete native source of the article as distributed in the arXiv e-print for this exact version is bundled unchanged as \texttt{sources/197732/source0.tex}.
\begin{equation}
T(t,p)=
\frac{(-1)^{p-1}}{2(p-1)!}\left(\psi^{(p)}(t)
-2\gamma\psi^{(p-1)}(t)-\sum_
{k=0}^{p-1} \binom{p-1}{k}\psi^{(k)}(t)\psi^{(p-1-k)}(t)
\right)\label{eq:Ttp}.
\end{equation}

\begin{theorem} \label{thm:thm3}
Let $C_{i,j}$ be the coefficients of the partial fraction decomposition
\begin{equation}
R(k) = \sum_{j=1}^r\sum_{i=1}^{p_j}\frac{C_{i,j}}{(k+t_j)^i},
\end{equation}
with distinct~$t_j$, $C_{p_j,j}\neq0$ and integers $p_j \geq
1$ whose sum is at least~2. Then
\begin{align}
\sum_{k=1}^\infty R(k){H_k}=\sum_{j=1}^r\sum_{i=1}^{p_j}C_{i,j}T(t_j,i),
\label{form:theorem4a}
\end{align}
with $T$ the function from \cref{eq:Ttp}.
\end{theorem}

\begin{corollary} \label{cor:cor2}
For integers $m, n \geq 1, \gcd (m, n) = 1, q \geq 1$ and $p \geq q+2$, and for the $T$ notation as defined in \cref{eq:Ttp},
\[\sum_{k=1}^\infty \frac{k^q H_k}{(mk+n)^p} =\frac1{m^p} \sum_
{j=0}^q 
(-1)^j
\binom{q}{j} \left(\frac nm\right)^j T \left(\frac nm,p-q+j\right).\]
\end{corollary}

\begin{proof} By \cref{thm:thm3}, it is sufficient to compute the
coefficients of the partial fraction decomposition of $k^q/(mk+n)^p$
with respect to~$k$. The result then follows from
\[\frac{k^q}{(mk+n)^p}=\frac{m^{-p}k^q}{(k+n/m)^p}=\frac{m^{-p}}{(k+n/m)^
{p-q}}\left (1- \frac{n/m}
{k+n/m}\right)^q=m^{-p}\sum_{j=0}^q (-1)^j \binom{q}{j} \frac{(n/m)^j}
{
(k+n/m)^
{p-q+j}}.\qedhere\]
\end{proof}

\end{document}
